\section{Recommendations}

The move a leader can make on Monday costs no procurement. It is to stand up the
cadence as a process before it is a product. Concretely: publish autonomy tiers
for the workflow configurations already running agents, from a sandbox tier up
through a trusted tier; define the trailing window and the outcome signals each
tier is scored on; set the relegation triggers and the clock, borrowing the
standard's own four-week shape as a starting point; and default every new or
sparsely-observed workflow to the sandbox. This is a review rhythm a platform or
developer-experience team can run in a spreadsheet on the first day and encode in
gateway policy as it matures. None of it requires buying anything.

The cadence adjusts trust over time; it presupposes a gate that screens every
change regardless of the workflow's record, and that gate is the cheapest control
to stand up. Turn on pre-merge security scanning, blast-radius and impact gating,
and agentic code review on every pull request, agent-authored or not. Agentic code
review is the fastest of the three to roll out org-wide and the lowest-friction:
the same industry benchmark found that adding it lifted the share of pull requests
merged within thirty days by up to about five percentage points, and that it
recovers much of the yield agents lose when they open changes no human owns
\citep{aibench2026}. This is the preventive layer earned autonomy sits on; it is
not an alternative to the cadence.

The capability to build alongside it is the delivery evidence layer, and it
belongs in the same budget cycle as the MCP boundary work, not after it. The
boundary makes actions enforceable; the layer makes outcomes legible; the cadence
needs both. Sequencing the enforcement work first and the evidence work later
produces the exact state this paper opened on: a fleet whose actions are
perfectly recorded and whose value is unattributable.

When evaluating any system, internal or purchased, that claims to govern agent
autonomy, five questions separate an outcome-conditioned tool from an
activity-graded one:

\begin{itemize}[leftmargin=1.4em,itemsep=0.25em,topsep=0.3em]
  \item \textbf{Attribute at the change level.} Does it join governed actions to
  the commit, review, and deploy they produced, or does it stop at the action
  record?
  \item \textbf{Close the governed-action-to-shipped-result loop.} Can a granted
  scope be traced to the outcome it caused, so authority and result meet in one
  place?
  \item \textbf{Segment authorship.} Does it tell agent-authored change apart
  from human-authored change, so the attribution gap is closed rather
  than aggregated over?
  \item \textbf{Resist adoption proxies.} Does it condition trust on
  held-in-production outcomes and refuse activity and adoption counts, such as
  action volume or acceptance rate, as stand-ins for delivery?
  \item \textbf{Serve context over governed rails.} Does it return evidence to
  the agents over the same typed, cache-contracted, per-call-authorized transport
  the actions ran on?
\end{itemize}

The 2026-07-28 spec settles the question of whether an organization can see and
control what its agents do, action by action. It leaves open the question this
paper is about: whether the organization can tell what those actions delivered,
and condition trust on the answer. The substrate for that second question now
exists. The layer above it, and the cadence above that, are the work in front of
every team running agents at scale.
